\documentclass[10pt,a4paper]{amsart}
\usepackage[margin=19mm]{geometry}
\usepackage{amsmath,amssymb,mathtools}
\usepackage[scr=rsfs,cal=euler]{mathalfa}
\usepackage{xfrac}
\usepackage[unicode,hidelinks,bookmarks=false]{hyperref}
\newcommand{\ie}{i.e.\ }
\newcommand{\cf}{cf.\ }
\newcommand{\resp}{respectively\ }
\newcommand{\Subsection}{\subsection}
\providecommand{\nobreakdash}{\nobreak\hbox{-}}
\setlength{\emergencystretch}{2em}
\multlinegap0pt
\usepackage{xcolor}
\usepackage{bm}
\usepackage{tikz}
\usepackage[normalem]{ulem}
\newtheorem{lemma}{Lemma}
\newtheorem{corollary}{Corollary}
\title{Fault-tolerant syndrome extraction in \texorpdfstring{$[[n,1,3]]$}{[[n,1,3]]} non-CSS code family generated using measurements on graph states}
\author{Harsh Gupta, Mainak Bhattacharyya, Ritik Jain, Ankur Raina}
\date{Source: arXiv:2501.12072v3 (CC BY 4.0)}
\begin{document}
\maketitle

\noindent\emph{Source and licence.} Fault-tolerant syndrome extraction in \texorpdfstring{$[[n,1,3]]$}{[[n,1,3]]} non-CSS code family generated using measurements on graph states. Version arXiv:2501.12072v3 (CC BY 4.0), distributed by arXiv under the Creative Commons Attribution 4.0 International licence, http://creativecommons.org/licenses/by/4.0/, as recorded in the arXiv licence metadata for this exact version and shown on the abstract page https://arxiv.org/abs/2501.12072v3. Full licence notice: \texttt{licenses/arxiv-2501.12072-CC-BY-4.0.txt}.\par\smallskip

\noindent\textit{Editorial scope.} Counted unit: the article's corollary 1 (source0.tex lines 538--546). The article's complete original proof of it is delivered with it (source0.tex lines 547--586, one \texttt{proof} environment of 40 lines). No mathematics is written, repaired or re-derived: the statement and the proof are the article's own lines, in the article's own notation, and no retained line is substituted or re-flowed.  Retained uncounted as the source-local context the proof needs: lines 474--477; lines 518--524.  Nothing was omitted from any retained span, so the package carries no omission record.  No retained line contains a citation, so the wrapper carries no bibliography and no bibliography entry.  No retained line contains a figure environment, an image-inclusion command or drawing code, so no image is delivered and the package declares no asset dependency.  Editorial formatting only, in addition: an AMS \texttt{amsart} preamble that restates the article's own declarations and macros used by the retained lines, namely the environments \texttt{lemma}, \texttt{corollary} on separate counters, together with the standard packages those lines need.  The retained environments print in this document as Lemma 1, Corollary 1 in order of appearance, whereas the article prints the same items with the numbering of its own full document.  The complete native source of the article as distributed in the arXiv e-print for this exact version is bundled unchanged as \texttt{sources/171901/source0.tex}.
\begin{equation}
 \mathrm{S}_u \;=\; 2^{\,n-k}-1 \;-\; \kappa(\mathbf{H}_{\mathrm{xyz}}).
 \label{eq:unique_syndrome}
\end{equation}

\begin{lemma}
\label{lemma:one qubit error}
For any stabilizer generator $g \in \mathcal{S} : g= \prod_{i} P_{a_i}$ of the $[[n,1,3]]$ stabilizer QECCs, if the weight of the stabilizer satisfies $w_g > 3$, then the number of uncorrectable hook errors is given by
\begin{equation*}
|\mathcal{U}_g|\leq w_g-3.
\end{equation*}
\end{lemma}

\begin{corollary}
\label{corollary: 1}

For a given $[[n,1,3]]$ code, the total number of abundant syndromes that can be assigned to the uncorrectable hook errors must satisfy
\[
\mathrm{S}_u \;\ge\; \sum_{\substack{g \in \mathcal{S}\\ w_g > 3}} \big(w_g - 3\big),
\]
where $\mathrm{S}_{u}$ is defined in \eqref{eq:unique_syndrome} and $w_g$ is the weight of the stabilizer $g$.
\end{corollary}

\begin{proof}















From Lemma \ref{lemma:one qubit error}, the number of uncorrectable hook errors generated during the measurement of a stabilizer generator \(g\) satisfies
$|\mathcal{U}_g|\leq w_g-3 $.
Let $$\mathcal{U}
:=
\bigcup_{\substack{g\in\mathcal{S}\\ w_g>3}}
\mathcal{U}_g$$
denote the set of all uncorrectable hook errors generated by the
stabilizer measurements. Then,





% % In Lemma \ref{lemma:one qubit error}, we proved that the number of uncorrectable hook errors for a given stabilizer $g$ of the $[[n,1,3]]$ code satisfies $|\mathcal U_g|\le w_g -3, \forall g \in \mathcal{S}$. Hence 
% \textcolor{red}{Define $\mathcal{U}$}
\[
|\mathcal U| \le \sum_g|\mathcal U_g| = \sum_{\substack{g\in\mathcal S\\w_g>3}}(w_g-3),
\]
% where  $\mathcal{U}$ is the set of all hook errors required for bare code.
Therefore, if the consequent bare code is fault-tolerant against all the hook errors, then it must contain a set of unique syndromes, cardinality of which is bounded by
\[
\mathrm{S}_u \ge\sum_{\substack{g\in\mathcal S\\w_g>3}}(w_g-3)\ .
\]
\end{proof}

\end{document}
